Ineens was er regen.
Dat viel Emmy tegen.
Ze was gehard tegen regen,

wind als wel kou,
maar naakt deed het toch au.
Rood en vervolgens blauw

kleurde haar onderlijf.
Kilheid stond buiten kijf.
Haar tepels werden stijf.

Ze wilde niet miepen.
Toen haar handjes sliepen,
haar vingers volliepen,

afgeknepen door touw,
vechtend tegen de kou,
bad ze aan God maar gauw

wat haar zo had gegrieft:
“Stad kom snel alsjeblieft!
Deze vrieskou doorklieft

mijn ontblote lichaam.
Ik onderkoel langzaam.”
Maar ze bleef gehoorzaam.

De regen werd onweer.
Emmy voelde afweer.
Weerstand tegen het weer.

Het was een voorbode,
God zijn alarmcode,
voor een periode

van onrecht en onheil.
Daardoor in allerijl
gingen ze onder zeil.

Toen het duister inviel,
werd de slavin labiel
en verbrak ze de deal.

Emmy nam de benen,
maar viel over stenen
en stootte haar schenen.

De knielende klager
smeekte slaafs de jager,
haar ijskoude plager,

om zijn lichaamswarmte.
Hij behield zijn kalmte.
Als volleerd beambte

had hij al in zijn hand
een schaapsleren halsband.
Die ging aan een leiband.

Vlot waarna een O-ring
om Emmy haar nek ging.
De ijzeren ketting

hield hij strak in zijn hand.
Hij trok haar naar zijn kant,
zodat haar mond beland

net onder zijn broekriem.
Ze duldde zijn regime
en stilde hem intiem.

De jager predikte
terwijl zij hem likte
of zijn eikel slikte.

“Luister goed naar mijn wet.
Je bent als slaafse slet
op de aarde gezet

om mij te verwennen.
Dat moet je erkennen.
Voor dit lot wegrennen

gaat in tegen Gods wil.
Je houdt je vooral stil.
Je gaat met mij van bil

zo vaak als ik verlang.
Je dient ook zonder dwang
het kerkelijk belang.

Je blijft mijn eigendom.
Waarin slechts ik klaar kom.
Behalve achterom.

Zie jezelf als puppy.
Altijd in je nakie.
Altijd vol met passie.

Altijd dol op spenen.
Ik zal je gaan trainen
om me bij te benen.

Ik zal je verpachten
en laten verkrachten,
lopende mijn wachten

op een hernieuwde zin.
Zo breng je mijn geld in.
Je bent mijn seksslavin.

Ik bepaal wat je doet
en wie je zwoel ontmoet.
Je verzorgt je zo goed

dat ik je lekker vind.
Zo word je vaak bemind.
Zo schenk je me een kind.

Als je mijn wet beloofd
en je daarvoor uitslooft,
buig dan knikkend je hoofd,

dan sta ik je niet af
voor je wettige straf.”
Natte Emmy stond paf.

De jonge hertenkalf,
slikkend zijn witte zalf,
hoorde bezig maar half

het groteske voorstel
tijdens het minnespel.
In ieder geval wel

zijn sterke kinderwens.
Ze miste de tendens.
Ze bevoelde haar pens

en werd alsmaar banger.
Was ze nu al zwanger
of duurde dat langer?

Ze had al wel gepaard.
Ze wilde geen bastaard.
Dat was haar alles waard!

Emmy’s brein ontplofte.
Ze zag zijn gelofte
als een trouwbelofte.

Wat een creatieve
draai door de naïeve.
Deze cognitieve

dissonantie bracht haar,
met vers zaad in haar haar,
tot een knikkend gebaar.

Vast in zijn matingnet
deed ze nog één schaakzet.
Zo had ze zich verwed.

Zo had ze daad bij woord
volbracht na dit akkoord,
dat ze deels had gehoord

en haar ten gronde richt.
Haar bespoten gezicht
persde de jager dicht

tegen zijn vlakke kruis.
Emmy haar nieuwe thuis.
Een bejaardentehuis,

haar tienerdagverblijf.
Zijn hete mannenlijf,
en het stroeve gewrijf

met eeltige handen,
die met ruwe randen
over haar huid brandden,

verdreven kort de kou.
Dat was wat ze graag wou.
Ze bleef zijn penis trouw.

Gelaatzaad

Mahabarata heeft een verhaal over een koningin uit het water die trouwt met een koning als die beloofd haar nooit iets te weigeren of tegen te werken. Uit geilheid gaat hij akkoord tot uw ziet dat ze zijn vijf zonen verdrinkt naar hun geboorte . De les is mak niet gedachteloos een belofte ( in dit geval uit lust)

het verhaal dat je beschrijft, komt uit de **mahabharata**, een van de grootste en belangrijkste epen uit de indiase literatuur. specifiek maakt het deel uit van de **adiparva**, het eerste boek van de mahabharata, waarin veel achtergrondverhalen en mythen worden verteld die de context scheppen voor het hoofdverhaal.

de koning in dit verhaal is **santanu**, en de koningin uit het water is **ganga**, de personificatie van de heilige rivier de ganges. het verhaal staat bekend als het **verhaal van santanu en ganga** of **santanu-ganga samvada** (de dialoog tussen santanu en ganga).

\#\#\# hoe ernaar te verwijzen:
1. **in tekst:**
- "zoals in het verhaal van santanu en ganga uit de mahabharata, waar een ondoordachte belofte tot tragedie leidde..."
- "het lot van santanu, die zijn zonen verloor door zijn blinde belofte aan ganga, herinnert ons aan..."
- "in de adiparva van de mahabharata wordt verteld hoe ganga, de riviergeest, santanu confronteerde met de prijs van zijn begeerte..."

2. **in voetnoten of bronvermelding:**
- "verhaal van santanu en ganga, mahabharata, adiparva, hoofdstuk [x]."
- "zie: mahabharata, boek 1 (adiparva), het verhaal van santanu en ganga."

3. **in dialoog of monoloog:**
- "weet je nog het verhaal van santanu? hij verloor alles door één onbedachtzaam moment van verlangen."
- "ganga, de riviergeest, was geen vrouw om lichtvaardig een belofte aan te doen."

als je dit verhaal in emmy nativus wilt verwerken, kun je het gebruiken als een metafoor, een verwijzing, of zelfs als een spiegelverhaal voor een van je personages. laat me weten of je hulp nodig hebt bij het integreren ervan!
